\chapter{Software Infrastructure}

\section{Version Control System Subversion}

Subversion is a  free/open source version control system,  which manages files
and  directories, and  the changes  made to  them, over  time. This  allows to
recover older versions and to have different developing branches \cite{svn}.

\subsection{Using Git in Combination with SVN}

Git\footnote{Git home page: \url{http://git-scm.com}} is a distributed version
control  system.    It  is  an  alternative  to   the  centralized  Subversion
system. The main difference is, that there may also exist a central repository
for Git but the local working  copies of this central repository are clones of
it and are therefore repositories themselves. On these local repositories, all
operations  for which  an interaction  with  the central  repository would  be
necessary using SVN (e.g. committing changes, creating a branch, etc.)  may be
performed  locally. The  local  repository  has to  be  synchronized with  the
central one at some points  in time.  This philosophy is therefore well-suited
for  working on  laptops for  example. Fortunately  there exists  also  a link
between SVN  and Git  which allows cloning  of a  central SVN repository  to a
local   Git  repository.   Table   \ref{tab:git_materials}  contains   further
references concerning Git.

The following script shows a typical sequence of commands for obtaining a copy
of the  source code, updating the  working copy, committing  local changes and
committing local changes back to the SVN repository.

\input{pygmentized/git_checkout_cfs}

\begin{table}[tb]
\centering
\begin{minipage}{18cc}\caption{Git resources on the web.}
\label{tab:git_materials}
\end{minipage}
\par\medskip\noindent
\begin{tabular}{ll}
  \toprule
  {\bf Description} & {\bf URL} \\ \otoprule
  Howto use Git and svn together & \url{http://flavio.castelli.name/howto_use_git_with_svn} \\
  Git in combination with SVN & \url{http://git.or.cz/course/svn.html} \\
  Git/SVN Howtos on Amarok Devel Site & \url{http://amaroklive.com/wiki/Development/Git} \\
  Professional Git (Book) & \url{http://progit.org/book/} \\
  Git Magic (Book) & \url{http://www-cs-students.stanford.edu/ ~blynn/gitmagic/index.html} \\
  Git Reference & \url{http://gitref.org/index.html} \\
  Reference Sheet & \url{http://cheat.errtheblog.com/s/git} \\
  \bottomrule
\end{tabular} 
\end{table}

\section{CMake Build System}

CMake is  a family of platform  independent tools designed to  build, test and
package  software  \cite{cmake,  Schroeder2004,  Martin2007a}.  It  is  mainly
targeted towards C,  C++ or Fortran projects. Instead  of writing Makefiles or
autoconf/automake scripts  on Unix or  Visual Studio projects on  Windows, one
just  needs to maintain  platform independent  configuration files  from which
CMake will  generate the  previously mentioned projects  for the  native build
environments.

\section{CFSDEPS}

 When  developing a  general purpose  finite element  code a  lot  of external
libraries (e.g.   BLAS, LAPACK, ARPACK, Metis\ldots)  have to be  pulled in if
the wheel  is not to  be reinvented.  Many  of these libraries are  present on
modern-day Linux/Unix  distributions. If more  than the standard  GNU compiler
suite should be supported it  is however almost certain that incompatibilities
with prebuilt system libraries will arise.

 On  Windows, not  all  of the  libraries  may be  readily available.  Another
problem which arises if the Visual C++ compilers are used, is that there exist
a number  of different  C and C++  runtime environments (static  vs.  dynamic,
single-threaded  vs. multi-threaded,  debug vs.  release, etc...).   A library
compiled and linked  against one runtime environment, must  not be linked with
libraries built against a different one.  Moreover, 64-bit libraries are still
rare to find.

Due to all of the aforementioned  problems, a library may not be available and
it  has to  be  built from  source.  Therefore  we  decided to  build all  our
dependencies  in an  automated script  based  manner.  The  scripts fetch  the
sources from the web, unpack  them, perform checksum checks, patch, configure,
build and  install the libraries (cf.  Fig. \ref{fig:cfsdeps_flowdiag}).  This
way we obtain maximum platform-independence while having a well-documented and
SVN-versioned description of how to build external packages. The necessity for
automatically building external packages has also been recognized by the CMake
developers   in  the   meantime   and  is   supported   since  version   2.8.0
(cf. \cite{BuildExtPackages})

The   current  hierarchical   structure  of   the  CFSDEPS   is   depicted  in
Fig.  \ref{fig:cfsdeps_structure} while  the  sequence of  actions during  the
build of an individual package is depicted in Fig. \ref{fig:cfsdeps_flowdiag}.

\begin{figure}[htb] 
  \centering
  \begin{overpic}[width=1.0\textwidth]{pics/cfsdeps_structure}
%    \put(10,20){\textcolor{white}{\bf The electrostatic man!}}
  \end{overpic}
  \caption{Structure and supported platforms of CFSDEPS.}
  \label{fig:cfsdeps_structure}
\end{figure}%

\begin{figure}[htb] 
  \centering
  \begin{overpic}[width=0.55\textwidth]{pics/cfsdeps_flowdiag}
%    \put(10,20){\textcolor{white}{\bf The electrostatic man!}}
  \end{overpic}
  \caption{Actions taken during the build of a package.}
  \label{fig:cfsdeps_flowdiag}
\end{figure}%


\section{CTest}

In order to  keep the quality of  a software project high it  is imparative to
perform manual and automated tests.  The necessary tests can be specified also
in the same  simple CMake syntax.  Every user should run  those tests before a
commit of  changes or additional software code.  This way it can  be made sure
that all the old features are still working at least on the user's platform.

\textcolor{red}{PLEASE NOTE:  The user name  and password for the  Ubuntu 8.04
VBox on rom are 'user' and 'user'.}

\section{CDash}

This is an  open source, web-based software testing  server, which aggregates,
analyzes and displays the results of software testing processes submitted from
different clients \cite{cdash}.  We use it to gather information about nightly
builds  on various  platforms. This  way we  can track  down errors  which are
platform-dependent (e.g. 32 bit vs.  64 bit) more efficiently.


\section{Trac}

This  tool  is  an  enhanced  wiki  and issue  tracking  system  for  software
development projects. Therewith, we can discuss different issues, fix the next
development  steps   by  assigning   tickets,  browse  the   {\sf  Subversion}
repository, can look at the changes been made between different versions, etc.
\cite{trac}